% !TEX program = lualatex
\documentclass[11pt,a4paper]{article}

\usepackage{luatexja}
\usepackage{amsmath,amssymb,amsthm}
\usepackage{geometry}
\geometry{margin=1in}

%-------------------------------------------------
%  定理環境
%-------------------------------------------------
\theoremstyle{definition}
\newtheorem{definition}{定義}[section]
\newtheorem{axiom}{公理}[section]
\newtheorem{theorem}{定理}[section]
\newtheorem{corollary}{系}[section]

%-------------------------------------------------
%  独自マクロ
%-------------------------------------------------
\newcommand{\dfix}{D_{\mathrm{fix}0}} % 0次元不動点（空）
\newcommand{\mweq}{\mathrel{\overset{\mathrm{mw}}{=}}}
\newcommand{\meta}{\Uparrow}
\newcommand{\Kvoid}{K_{\emptyset}}
\newcommand{\Bou}{\mathbf{Bou}}
\newcommand{\nCat}{\mathbf{nCat}}
\newcommand{\Leadsto}{\rightsquigarrow}
\newcommand{\susp}[1]{#1^*}
\newcommand{\Ymw}{Y_{\mathrm{mw}}}

\renewcommand{\abstractname}{Abstract}

\begin{document}

\title{1Bou:nCat：プログラム実行システムのまとめ}
\author{Masaomi Chiba}
\date{2026-04-25}
\maketitle

\begin{abstract}
本稿は、界論（Bou）と下位圏（nCat）を統合した
1Bou:nCat のプログラム実行システムをまとめる。
世俗諦的な構造（nCat）と勝義諦的な構造（Bou）の間の
埋め込み・遷移・完備化・解脱のプロセスを形式化し、
任意の形式体系 $N$ に対し $N \otimes I_N$ が実行可能なプログラミング言語となる
ことを示す。
\end{abstract}

\section{1Bou:nCat の構成}

\begin{definition}[1Bou:nCat]
1Bou:nCat は、以下の二つの層から成る。
\begin{itemize}
\item $\nCat$：下位圏の層（世俗諦）  
  - 例：CCC（型付きラムダ計算）、モノイダル圏（線形論理）、LCCC（依存型）、確率的圏、線形圏など
\item $\Bou$：界論の層（勝義諦）  
  - 演算子：$\Leadsto$（遷移）、$\meta$（メタ否定）、$\susp{}$（保留）、$\dfix$（絶対不動点）
\end{itemize}
\end{definition}

\begin{definition}[埋め込み関手 $\mathcal{E}$]
下位圏 $\nCat$ から界 $\Bou$ への埋め込み関手 $\mathcal{E}$ を定義する。
\begin{align*}
\mathcal{E} &: \nCat \hookrightarrow \Bou \\
\mathcal{E}(A) &= \susp{A} \quad\text{（対象の保留）} \\
\mathcal{E}(f) &= A \meta B \quad\text{（射のメタ否定）}
\end{align*}
\end{definition}

\section{プログラム実行のプロセス}

\begin{definition}[帰納フラクタロイド $I_N$]
古典的な「モナド」は自己同一性と平坦な結合律を要求するため、本体系には適合しない。
代わって、界 $\Bou$ 上で動的プロセスを生成する文脈演算子を
「帰納フラクタロイド（Inductive Fractaloid）$I_N$」と定義する。
$I_N$ は、中道不動点コンビネータ $\Ymw$ を駆動エンジンとし、
静的な対象を「保留・遷移・解脱」のフラクタル軌道へと乗せる作用素である。
\end{definition}

\begin{definition}[プログラム $P$ の実行結合 $\otimes$]
任意の形式体系 $N$（論理・文法・計算：$\nCat$ 上の構造）に対し、
プログラム $P$ は、静的構造 $N$ と動的フラクタロイド $I_N$ の
結合（着火）として定義される。
\begin{align*}
P &:= N \otimes I_N \\
\text{実行軌道} &: \susp{P_0} \Leadsto P_1 \Leadsto \dots \Leadsto \meta^4 P_k \mweq \dfix
\end{align*}
静的対象は $\Bou$ に埋め込まれた瞬間に保留（$\susp{P_0}$）され、
$\Ymw$ の作用により状態遷移（$\Leadsto$）を繰り返し、最終的に $\dfix$ へと解脱する。
\end{definition}

\begin{theorem}[$N \otimes I_N$ の実行可能性と停止性]
任意の形式体系 $N$ に対し、$N \otimes I_N$ は実行可能なプログラミング言語となる。
\end{theorem}
\begin{proof}
$N$ を埋め込み関手 $\mathcal{E}$ により界 $\Bou$ へ移行させる。
$\Bou$ 上において、自己同一性を要求する古典的評価は行われず、
帰納フラクタロイド $I_N$ に従った中道推論（大域的カット消去）が適用される。
$\Ymw$ による自己参照（無明の生成）は、系に蓄積された構造的負債（Cut）を
消費しながら進行し、負債が尽きた時点で $\meta^4$ の相転移を起こし、
無記状態 $\dfix$ へと必然的に収束（停止）する。
\end{proof}

\section{典型例}

\begin{corollary}[ペアノ算術 $\mathbb{P}$]
$\mathbb{P} \otimes \text{Free}(\Sigma_+)$ は項書換え系・純関数言語となる。
\end{corollary}

\begin{corollary}[最小論理 $\mathrm{ML}$]
$\mathrm{ML} \otimes \text{Cont}_R$ は古典論理 + call/cc となる。
\end{corollary}

\begin{corollary}[文脈自由文法 $\mathrm{CFG}$]
$\mathrm{CFG} \otimes \text{State}_{\mathbb{N}}$ は再帰下降パーザとなる。
\end{corollary}

\begin{corollary}[ホーン節 $\mathrm{Horn}$]
$\mathrm{Horn} \otimes \text{List}$ は Prolog となる。
\end{corollary}

\section{界論（Bou）の演算子とプログラム実行}

\begin{definition}[遷移 $\Leadsto$]
プログラム $P$ の状態遷移を表す。
$$P \Leadsto P'$$
は、$P$ が次の状態 $P'$ へ遷移することを意味する。
\end{definition}

\begin{definition}[メタ否定 $\meta$]
世俗諦から勝義諦へのメタレベル上昇を表す。
$$A \meta B$$
は、$A$ の境界を生成し、知覚次元を一段引き上げる。
\end{definition}

\begin{definition}[保留 $\susp{}$]
局所的境界を生成し、世俗諦へ戻る操作を表す。
$$\susp{A}$$
は、$A$ を保留状態に置き、局所的な境界を生成する。
\end{definition}

\begin{definition}[絶対不動点 $\dfix$]
すべての境界が融解した極限状態を表す。
$$\meta^4 A \mweq \dfix$$
プログラム実行の終了点・解脱点として機能する。
\end{definition}

\section{中道不動点コンビネータ $\Ymw$}

\begin{definition}[中道不動点コンビネータ $\Ymw$]
$\meta$ に対し、以下の中道等価を満たすコンビネータ $\Ymw$ を定義する。
$$\Ymw\, \meta \mweq \meta\,(\Ymw\, \meta)$$
$\Ymw$ はフラクタル四則から必然的に導出される
「無明（自己参照の閉じ損ね）」のジェネレータであり、
$\Kvoid$ への相転移（大域的カット消去）によって停止する。
\end{definition}

\section*{備考：1Bou:nCat の意義}

1Bou:nCat は、
- 世俗諦的な構造（nCat）と勝義諦的な構造（Bou）を統合し、
- 任意の形式体系 $N$ に対し $N \otimes I_N$ が実行可能なプログラミング言語となる
ことを示した。

これにより、
- 論理・文法・計算・確率・線形性など、多様な構造を
- 一つの枠組み（1Bou:nCat）で統一的に扱える
という理論的基盤が確立された。

\end{document}
